\chapter{Considerações Finais}
\label{cap:consideracoes}

Este trabalho perguntou quanto os questionários contextuais de professor e
de diretor acrescentam, fora da amostra, à previsão dos desfechos das
escolas públicas do 2º ano além do INSE e da geografia. A resposta é que
acrescentam pouco. Na proficiência média, os dois questionários juntos
reduzem o erro em 4,2\%, com intervalo de confiança que exclui zero. Na
proporção de alfabetizados, o ganho não se distingue de zero. Das quatro
hipóteses, só a do INSE (H3) é confirmada nas duas trilhas. A do diretor
(H2) fica perto da margem sem confirmá-la, e as do professor (H1) e da
geografia (H4) não se confirmam. O modelo linear regularizado foi
selecionado nas duas trilhas, porque os \textit{ensembles} não superaram a
regra de seleção definida antes da estimação.

O trabalho deixa quatro contribuições, as mesmas que a
Seção~\ref{sec:justificativa} anunciou. No plano metodológico, compara
modelos lineares e \textit{ensembles} em duas trilhas no nível da escola,
com o mesmo conjunto de preditores, e mostra que o INSE pesa nas duas
trilhas e que a ordem dos demais temas muda de uma para a outra.
No rigor de dados, registra a decisão de inclusão de cada variável
candidata, trata o vazamento do alvo e mantém cor/raça como variável de
auditoria, com as duas especificações lado a lado. Para a política pública,
mostra que o INSE e a geografia já entregam a maior parte da capacidade de
prever o desempenho das escolas e que o erro é maior nas escolas de INSE
baixo, o que desaconselha usar a previsão para classificar escolas
individuais. Para a Ciência de Dados para Negócios, percorre o ciclo
inteiro, da integração de quatro bases a um \textit{pipeline}
reprodutível, com seleção de modelo pré-registrada e explicações
auditáveis.

As limitações estão na Seção~\ref{sec:discussao}. As mais importantes são
o caráter ecológico e transversal da análise, a ausência de medida do
método de alfabetização no questionário do professor e o tamanho efetivo
pequeno das médias ponderadas.

Três linhas de trabalho futuro decorrem desses limites. A primeira é
estimar o efeito de uma prática modificável apontada pelo ranking de
temas, como as do bloco do diretor, com várias edições do Saeb ou com um
desenho quase experimental. Florestas causais \cite{wager2018causal} e
\textit{double machine learning} \cite{chernozhukov2018double}, com
implementações como a do CausalML \cite{chen2020causalml}, servem a esse
fim desde que as hipóteses de identificação discutidas na
Seção~\ref{sec:registros-inferenciais} sejam plausíveis. A segunda é testar
se os temas do diretor que mais pesam no modelo descrevem o tipo de escola,
cruzando-os com o porte, o tempo de permanência dos alunos e a rede. A
terceira depende do instrumento. Incluir no questionário do professor
perguntas sobre o método de alfabetização permitiria levar à escola a
evidência experimental sobre o ensino da leitura revisada na
Seção~\ref{sec:alfabetizacao-2ef}.
